\chapter*{Introducción}

\noindent\textbf{Situación Problemática}

\noindent La agricultura cubana enfrenta una crisis multiplicativa que involucra la escasez de materias primas, energía, agua y suelo, donde cada factor degradado amplifica el impacto de los demás. Esta situación se manifiesta en una crisis hídrica generalizada caracterizada por sequías prolongadas, acumulados de lluvia por debajo de los valores medios históricos y un patrón errático: escasez prolongada interrumpida por tormentas torrenciales aisladas que dañan físicamente los cultivos de plátano, maíz y frutales, además de causar erosión severa. El cambio climático exacerba estas condiciones mediante lluvias erráticas, olas de calor, huracanes e inundaciones que alteran la planificación agrícola \citep{cigg2025}.

Paralelamente, la infraestructura de riego se encuentra colapsada, lo cual constituye una contradicción crítica dado que se necesita precisamente para combatir la sequía. Una pequeña parte de los cultivos dispone de sistemas de riego, predominando tecnologías obsoletas como riego por aspersión, gravedad, aniego y la utilización de tracción animal, lo que provoca bajo rendimiento por hectárea \citep{cornelltraccion}. Existe ausencia de precisión agrícola: siembra a voleo sin dosificación precisa de semillas, fertilización uniforme sin análisis de suelo y cosecha manual en cultivos mecanizables. La gestión institucional resulta deficiente por falta de personal para operación y control de canales, programación de riego sin balances hídricos; se riega ``por calendario'' y no por necesidad real del cultivo.

La intersección fatal con las materias primas, mayormente provocada por el bloqueo económico y las medidas coercitivas, determina que sin combustible no haya electricidad para bombear agua de embalses ni pozos hacia los canales de riego. La maquinaria de bombeo está paralizada por falta de piezas y los apagones programados impiden el funcionamiento de sistemas de riego por goteo o aspersión más eficientes. A esto se suma la salinización de los suelos provocada por la incremental necesidad de producción agrícola para satisfacer la demanda, sobreexplotando suelos marginales sin fertilizantes, tecnología de riego de precisión y drenaje adecuado, lo que genera estrés salino haciendo inútil que se provea de agua.

Cuba es un Estado insular caribeño de 110\,860~km\textsuperscript{2} ubicado en la zona intertropical, con clima cálido húmedo (temperatura media 25\,°C) y alta vulnerabilidad a fenómenos meteorológicos extremos como huracanes, sequías y ciclones tropicales. Geográficamente, el país presenta una marcada estacionalidad en precipitaciones, con un período seco (noviembre-abril) que coincide con la campaña agrícola de frío, y un período lluvioso con riesgo de huracanes (mayo-octubre) \citep{onei2024}.

Históricamente, la agricultura cubana ha transitado por modelos productivos intensivos en insumos externos: primero el monocultivo azucarero colonial, luego la mecanización soviética (1960-1990) basada en subsidios energéticos, y finalmente el ``Período Especial'' de los años 90 que provocó un colapso tecnológico forzado y el retorno parcial a la tracción animal. Desde entonces, el sector opera bajo una economía estatal centralizada con reformas parciales, pero manteniendo dependencia crítica de importaciones de alimentos (70\,\% del consumo interno) e insumos agrícolas que el país no produce localmente \citep{rosset2020}. El marco institucional actual comprende el Ministerio de Agricultura (MINAG) y el Instituto Nacional de Recursos Hidráulicos (INRH), bajo políticas de soberanía alimentaria y seguridad nacional programadas hasta 2030; sin embargo, el sector opera en condiciones de asfixia financiera por el bloqueo comercial estadounidense y la crisis de divisas \citep{alnap2024, leyssan2022}.

\bigskip
\noindent\textbf{Problema}

\noindent ¿Cómo diseñar un sistema basado en tecnologías IoT que permita la gestión automática y eficiente del riego agrícola, optimizando el uso del agua en condiciones de limitaciones energéticas y tecnológicas propias del contexto cubano?

\bigskip
\noindent\textbf{Hipótesis}

\noindent Si se implementa un sistema IoT de bajo costo denominado FoT (\textit{Farm of Things}), basado en Arduino UNO y sensores ambientales, con capacidad de operación autónoma ante interrupciones eléctricas, entonces es posible automatizar la gestión del riego agrícola reduciendo el consumo de agua respecto a los métodos tradicionales y disminuyendo la dependencia de intervención humana constante.

\bigskip
\noindent\textbf{Objeto de Estudio}

\noindent Las herramientas de transmisión de datos en entornos de IoT aplicadas a la agricultura de precisión.

\bigskip
\noindent\textbf{Campo de Acción}

\noindent El diseño e implementación de sistemas para el monitoreo y control automatizado del riego en contextos de recursos limitados.

\bigskip
\noindent\textbf{Objetivo General}

\noindent Desarrollar un sistema IoT de bajo costo que permita la autogestión del riego agrícola mediante el monitoreo en tiempo real de variables ambientales del suelo, optimizando el uso de los recursos.

\bigskip
\noindent\textbf{Objetivos Específicos}

\begin{enumerate}
	\item Asimilar la tecnología IoT identificando las herramientas y componentes 
	ventajosos para el uso de Arduino UNO en el desarrollo del sistema FoT.
	\begin{enumerate}
		\item Estudiar los fundamentos teóricos del IoT aplicado a la agricultura 
		de precisión y revisar soluciones similares existentes.
		\item Seleccionar los componentes de hardware (Arduino UNO, sensores de 
		humedad, temperatura y conductividad) más adecuados al contexto.
		\item Identificar los protocolos de comunicación y mecanismos de 
		persistencia de datos viables bajo restricciones energéticas.
	\end{enumerate}
	
	\item Analizar el negocio para extraer los artefactos necesarios mediante la 
	identificación de requisitos, el diseño de la arquitectura propuesta y su despliegue.
	\begin{enumerate}
		\item Identificar y documentar los requisitos funcionales, atributos de 
		calidad y restricciones del sistema.
		\item Diseñar la arquitectura candidata del sistema FoT, incluyendo 
		diagramas de casos de uso y mecanismos de diseño.
		\item Implementar el prototipo del sistema integrando los módulos de 
		monitoreo, decisión y actuación sobre el riego.
	\end{enumerate}
	
	\item Validar el sistema de autogestión demostrando que opera de forma eficaz 
	con una disminución medible del consumo de recursos.
	\begin{enumerate}
		\item Definir las métricas e indicadores de evaluación: precisión de 
		sensores, autonomía ante cortes eléctricos y consumo de agua.
		\item Ejecutar pruebas de funcionamiento autónomo sobre un sembrado de 
		ejemplo bajo condiciones controladas.
		\item Comparar los resultados obtenidos con métodos de riego tradicionales 
		y verificar la reducción del consumo de recursos.
	\end{enumerate}
\end{enumerate}

\bigskip
\noindent\textbf{Aporte}

\noindent El aporte del presente trabajo radica en proporcionar una solución tecnológica adaptada a las restricciones económicas y energéticas del contexto cubano, demostrando la viabilidad de implementar agricultura de precisión con tecnologías de bajo costo. La factibilidad técnica se sustenta en componentes accesibles y documentación extensa; la factibilidad económica en costos relativamente bajos que permiten adquisición gradual, reduciendo la dependencia de grandes inversiones; y la factibilidad operativa en un diseño de bajo consumo energético adaptable a fuentes alternativas como baterías o paneles solares.

\bigskip
\noindent\textbf{Preguntas de Investigación}

\noindent Las preguntas de investigación que guían este trabajo son: ¿puede mantenerse activo el sistema a pesar de cortes bruscos de electricidad que dejen inactivo el sistema?; ¿es viable técnica y económicamente el mantenimiento de un sistema IoT agrícola en el contexto cubano?; ¿en qué medida el sistema propuesto reduce el consumo de agua en comparación con métodos tradicionales?; y ¿qué nivel de autonomía puede alcanzar el sistema sin intervención humana constante?

\bigskip
\noindent\textbf{Validación}

\noindent La validación se realizará mediante un pequeño sembrado de ejemplo para evaluación física, empleando herramientas externas para comparar con los sensores y mecanismos de medición con el sistema de autogestión. Los puntos a evaluar comprenden el análisis de precisión en la lectura de variables ambientales, la verificación del funcionamiento autónomo del sistema durante ciclos de riego, y la medición de la estabilidad del sistema frente a interrupciones eléctricas.

\bigskip
\noindent El presente trabajo se estructura en tres capítulos. El \textbf{Capítulo 1} presenta el marco teórico referencial, abarcando los fundamentos del Internet de las Cosas (IoT) aplicado a la agricultura, los sistemas de monitoreo de variables ambientales, la transmisión de datos en tiempo real bajo recursos limitados, y el análisis de soluciones similares. El \textbf{Capítulo 2} expone la solución propuesta: la documentación de ingeniería de requisitos, la arquitectura candidata del sistema FoT y los mecanismos de diseño identificados. El \textbf{Capítulo 3} describe la validación de la solución mediante pruebas físicas sobre un sembrado de ejemplo. Finalmente se presentan las conclusiones y recomendaciones derivadas de la investigación.